% XeLaTeXでコンパイルしてください．
\documentclass[a4paper,11pt]{article}
\usepackage{xeCJK}
\setCJKmainfont{HaranoAjiMincho-Regular.otf}[BoldFont=HaranoAjiMincho-Bold.otf]
\renewcommand{\thesection}{\arabic{section}}
\usepackage[margin=25mm]{geometry}
\usepackage{tikz}
\usetikzlibrary{calc}
\usepackage{enumitem}
\usepackage{amsmath}
\usepackage{booktabs}
\usepackage[hidelinks]{hyperref}
\renewcommand{\refname}{参考文献}

% 表紙の記載内容
\newcommand{\reporttitle}{アジャイルワーク２ 実験報告書}
\newcommand{\studentnumber}{25G1074}
\newcommand{\studentname}{泉水 悠斗}

\usepackage{tabularx,array,longtable,listings,pgfplots,caption}
\usepackage{xurl}
\usepackage{float}
\setCJKmonofont{HaranoAjiMincho-Regular.otf}
\Urlmuskip=0mu plus 1mu
\raggedbottom
\clubpenalty=10000
\widowpenalty=10000
\displaywidowpenalty=10000
\pgfplotsset{compat=1.18}
\usepackage{titlesec}
\titleformat{\section}{\large\bfseries}{\thesection}{1em}{}
\titleformat{\subsection}{\normalsize\bfseries}{\thesubsection}{1em}{}
\setlength{\parindent}{1em}
\setlength{\parskip}{3pt}
\renewcommand{\baselinestretch}{1.12}
\renewcommand{\figurename}{図}
\renewcommand{\tablename}{表}
\captionsetup{font=small,labelsep=quad}
\lstset{basicstyle=\ttfamily\footnotesize,columns=fullflexible,keepspaces=true,frame=single,breaklines=true,aboveskip=8pt,belowskip=8pt}
\newcolumntype{Y}{>{\raggedright\arraybackslash}X}
\newcommand{\code}[1]{\nolinkurl{#1}}
\begin{document}
\begin{titlepage}
\thispagestyle{empty}
\begin{tikzpicture}[remember picture,overlay]
    % A4用紙左上を原点として配置．画像の枠・影・余白を再現．
    \coordinate (origin) at (current page.north west);
    \fill[black] ($(origin)+(17.5mm,-33.5mm)$) rectangle ($(origin)+(202.5mm,-70mm)$);
    \draw[fill=white,line width=.4pt] ($(origin)+(15.5mm,-31.5mm)$) rectangle ($(origin)+(200.5mm,-68mm)$);
    \node[font=\bfseries\fontsize{25}{32}\selectfont] at ($(origin)+(108mm,-50mm)$) {\reporttitle};
    \node[align=center,font=\fontsize{17}{25}\selectfont] at ($(origin)+(108mm,-109mm)$) {千葉工業大学\\情報変革科学部 情報工学科};
    \draw[rounded corners=4mm,line width=.4pt] ($(origin)+(15.5mm,-200mm)$) rectangle ($(origin)+(202.5mm,-268mm)$);
    \node[anchor=west,font=\bfseries\fontsize{17}{22}\selectfont] at ($(origin)+(43mm,-226mm)$) {学生番号};
    \node[font=\bfseries\fontsize{17}{22}\selectfont] at ($(origin)+(121mm,-226mm)$) {\studentnumber};
    \draw ($(origin)+(73mm,-231mm)$) -- ($(origin)+(166mm,-231mm)$);
    \node[anchor=west,font=\bfseries\fontsize{17}{22}\selectfont] at ($(origin)+(43mm,-259mm)$) {氏\hspace{15mm}名};
    \node[font=\bfseries\fontsize{17}{22}\selectfont] at ($(origin)+(121mm,-259mm)$) {\studentname};
    \draw ($(origin)+(75mm,-264mm)$) -- ($(origin)+(168mm,-264mm)$);
\end{tikzpicture}
\null
\end{titlepage}


\setcounter{page}{1}
\section{実験の目的}
ArduinoによるRGB補正では，学習済みニューラルネットワークの重みと推論時の作業領域を限られたメモリに収める必要がある．
隠れ層を大きくすると表現できる関数の範囲は広がるが，重み数も増え，搭載可能なモデルの大きさが制約される．本実験では，枝刈りと疎行列保存によって，非圧縮では搭載できない大きさのモデルをArduino上で動作させられるかを調べる．
実験計画書\cite{plan}の目的を引き継ぎ，搭載できる隠れ層次元数，RGB補正画像の平均二乗誤差（MSE），
全パラメータ保存容量，平均推論時間を評価する．保存形式はCOOとCSRを比較し，
同一次元数の比較では同じ学習済み重みと枝刈り閾値を用いる．次元数の増加による精度の変化と，
形式だけを変更した際の容量・時間の差を分けて検討する．
本報告書は，2026年10月から11月6日までの保存済み測定結果を2026年11月13日に整理したものである．
報告書作成に伴う再学習・実機での再測定は行っていない．計画では第1隠れ層を40に固定していたが，
実施時は両隠れ層を同時に拡大した．変更の内容と理由は第\ref{sec:change}節に示す．

\section{実験の理論}
\subsection{モデルと学習}
実施した主なモデル構成は$3\to n\to n\to3$である．
入力を255で正規化したRGBベクトル$\boldsymbol{x}$とし，

\begin{align}
    \boldsymbol{h}_1&=\operatorname{ReLU}(W_1\boldsymbol{x}+\boldsymbol{b}_1),\\
    \boldsymbol{h}_2&=\operatorname{ReLU}(W_2\boldsymbol{h}_1+\boldsymbol{b}_2),\\
    \hat{\boldsymbol{y}}&=W_3\boldsymbol{h}_2+\boldsymbol{b}_3
\end{align}

で補正RGBを得る．$W_1$は$n\times3$，$W_2$は$n\times n$，$W_3$は$3\times n$であり，
出力層にはReLUを使用しない．100次元以上の学習では，正規化RGBのMSEに全3層の重みのL1ノルムを加え，

\[
    L=\frac{1}{3N}\sum_{p=1}^{N}\|\hat{\boldsymbol{y}}_p-\boldsymbol{y}_p\|_2^2
    +10^{-6}\sum_{\ell=1}^{3}\|W_\ell\|_1
\]

を最小化する．L1正則化は小さい重みを生じやすくするが，学習だけで厳密な0になることを保証しないため，
別途閾値による枝刈りを行う．

\subsection{枝刈りとCOO・CSR}
本実験では$W_1,W_2$に対し，共通の閾値$t$を用いて

\[
    w'_{ij}=\begin{cases}0&(|w_{ij}|<t),\\w_{ij}&(|w_{ij}|\geq t)\end{cases}
\]

とする．閾値と等しい重みは残し，$W_3$と全バイアスは変更しない．枝刈りはモデルの関数を変える処理であり，
その後に同じ非ゼロ値をCOO・CSRへ保存する処理とは区別する．通常配列に0を並べるだけでは保存容量は減らない．
枝刈り後の再学習は行っていない．Hanら\cite{han}の手法には再学習も含まれるため，
本実験がその全手順を再現したものではないことに注意する．
COOは非ゼロ要素ごとに値・行番号・列番号を保存する\cite{coo}．CSRは行順に並べた値・列番号・行の開始位置配列を保存する\cite{csr}．
例えば

\[
    W'=\begin{pmatrix}0&0.5&0\\0&0&0\\-0.2&0&0.8\end{pmatrix}
\]

では，値$[0.5,-0.2,0.8]$と列番号$[1,0,2]$は共通であり，
COOは行番号$[0,2,2]$，CSRは開始位置$[0,1,1,3]$を持つ．CSRの第$r$行は，
開始位置$r$から次の開始位置の直前までを参照する．空行でも開始位置の要素は必要である．
両形式が同じ行列を表せば，数学的な推論結果は同じになる．ただし，浮動小数点の加算順が変わる一般的な実装では
微小差が生じ得る．本実装は行順・列順をそろえて変換し，復元一致も確認している．

\subsection{保存容量と圧縮率}
重み・バイアスを4 Bのfloatで保持する非圧縮版の容量は

\[
    S_{\mathrm{dense}}(n)=4(3n+n+n^2+n+3n+3)=4(n^2+8n+3)\quad[\mathrm{B}]
\]

である．COO・CSRでは値に加えて索引の容量も必要になる．ある層の非ゼロ数を$K$，行数を$R$，
行索引・列索引・開始位置1要素のバイト数をそれぞれ$b_r,b_c,b_p$とすると，その重みの保存容量は

\[
    S_{\mathrm{COO}}=(4+b_r+b_c)K,\qquad
    S_{\mathrm{CSR}}=(4+b_c)K+b_p(R+1)
\]

となる．したがってCSRが小さくなる条件は$b_rK>b_p(R+1)$であり，極端に非ゼロ要素が少ないとCOOが小さくなり得る．
本実装は索引の最大値に応じて1，2，4 Bの符号なし整数を選ぶため，型の変化も容量に影響する．
非ゼロ数が0の配列にはダミー1要素を保持し，その容量も含める．
本報告書では，圧縮率を

\[
    C=\frac{\text{圧縮後の全パラメータ容量}}{\text{同一次元数の非圧縮全パラメータ容量}}\times100\quad[\%]
\]

とする．$C$は残った容量の割合であり，削減率は$100-C$である．ヘッダーの文字数，プログラム全体のフラッシュ量，静的RAM量とは区別する．

\section{計画からの変更とその理由}\label{sec:change}
\begin{table}[H]
\centering\small
\caption{計画と実施方法の主な相違}
\renewcommand{\arraystretch}{1.25}
\begin{tabularx}{\linewidth}{p{23mm}YY}\toprule
項目&計画書&実施した方法\\\midrule
モデル構成&$3\to40\to n\to3$&$3\to n\to n\to3$を主系列とした\\\midrule
圧縮対象&第2層の重みのみ&第1・第2層の重み．第3層と全バイアスは保持\\\midrule
次元数の刻み&10ずつ増加&非圧縮40，50，60，70．圧縮70，100と100刻みの条件を検証\\\midrule
枝刈り閾値&搭載可否とMSEを見て調整&70・100・700で0.005，0.01，0.02，0.03．200〜900は0.02を主条件とした\\\midrule
試験の終了&両形式の限界まで検証&800・900および700の0.005はCOOのRAM不足で停止．CSRの実機試験は未実施\\\bottomrule
\end{tabularx}
\end{table}

\subsection{両隠れ層を同時に拡大した理由}
計画の$3\to40\to n\to3$では，$W_2$の要素数は$40n$であり，
第1段で作る特徴数を固定したまま第2段だけを広げる実験となる．実施時は，
幅をそろえた2隠れ層モデル全体の規模を変え，規模拡大に対する搭載可否と補正性能を調べる構成に変更した．
これにより$W_2$は$n^2$要素となり，モデル拡大に伴う容量増加と疎化による削減効果を同じ構成系列で評価できる．
例えば$n=700$では，計画の非圧縮容量は$652+176\times700=123,852$ B，実施構成では1,982,412 Bとなる．
これは単に実装を容易にする変更ではなく，両段の特徴数を増やしたモデルの搭載範囲を対象にした変更である．
ただし，第2隠れ層だけを拡大した効果を切り分けることはできない．初期に行った40固定系列とは混同せず，
本報告の主な比較は両隠れ層が同じ次元数の系列に限定する．

\subsection{第1層も圧縮した理由}
両隠れ層を拡大すると，第1層の重みも$3n$個に増加する．第2層だけを圧縮しても，
第1層の保存・演算コストが残るため，第1層と第2層を同じ基準で疎化し，モデル全体としての搭載可能性を評価することにした．
第1層は列数が3と小さいため，索引を含めると必ず小さくなるわけではない．そこで，各層の非ゼロ数と索引型を含む全パラメータ容量で有効性を確認する．
700次元・閾値0.02では，第1層の非ゼロ数は47／2,100であり，COOの保存量は$47(4+2+1)=329$ Bとなる．
非圧縮の8,400 Bより8,071 B小さい．第1層の圧縮にも実際に容量削減があることが確認できる．
第3層は出力3成分を直接生成するため，まずその重みとバイアスを保持した条件で評価対象を絞った．
この選択が精度保持を保証するわけではなく，第3層の圧縮効果は今後の検証事項である．

\subsection{刻み幅・閾値と終了条件}
非圧縮の限界付近では10刻みを用い，圧縮後は100刻みを主として，広いモデル規模で容量・時間・MSEの傾向を調べた．
刻みを粗くすると検証条件を広げられる一方，限界次元数の厳密な特定はできない．700動作・800失敗から，
701〜799も動作不能と結論することはできない．また，実際の測定順には900，800，700の順で限界付近を確認した箇所があり，
単調に増加させただけではない．
閾値は0.01を基準に，より弱い枝刈り0.005と，より強い0.02，0.03を比較した．この4値は精度と容量の変化を見るための探索点であり，
最適値を理論的に保証するものではない．
先行する100次元の閾値比較では，0.02のMSEは83.370となり，検証した4閾値の中で最小であった．
0.01のMSEは91.778であり，0.02ではMSEに加え，COO・CSRともに保存容量と推論時間も小さくなった．
一方，0.03まで閾値を上げると保存容量と推論時間はさらに減少したものの，MSEは88.481へ増加した．
この結果から，補正精度を優先しながら小容量化も図るという本実験の目的に照らすと，共通閾値として用いた0.02は，
100次元の比較結果からも妥当と考えられる．ただし，この評価は試した4閾値と固定9画素に基づくものであり，
他の次元数や評価画像でも0.02が最良となることを保証するものではない．
200〜900の主系列は0.02にそろえ，次元数と閾値の影響が同時に変わらないようにした．
COOがメモリ不足になった時点で停止した条件では，未実施のCSRを「失敗」とは扱わない．

\section{実験環境と方法}
\subsection{使用環境と対象データ}

\begin{table}[H]\centering\small
\caption{100〜900次元の主系列に共通する条件}
\renewcommand{\arraystretch}{1.15}
\begin{tabularx}{\linewidth}{p{35mm}Y}\toprule
項目&条件\\\midrule
ボード&Arduino UNO R4 WiFi，RA4M1，48 MHz，SRAM 32,768 B，フラッシュ262,144 B\cite{board}\\\midrule
ビルド条件&FQBN \code{arduino:renesas_uno:unor4wifi}，core 1.5.1．モデル配列はRAM配置のまま比較\\\midrule
学習環境&保存された学習記録：Python 3.13.14，PyTorch 2.14.1，NumPy 2.5.3，OpenCV 5.0.0．CPU 1スレッド\\\midrule
学習条件&33画素の全件を使用，Adam，学習率0.01，50,000回，L1係数$10^{-6}$，seed 0，決定的アルゴリズム有効\\\midrule
評価条件&保存済み3$\times$3の9画素，行優先・RGB順，入力を255で除算．基準画像は\code{reference_3x3.ppm}\\\midrule
時間測定&予備100回，1セット1,000回，5セット．全条件で同じ入力順と集計処理\\\bottomrule
\end{tabularx}\end{table}

学習入力は\code{scanner_ver2/before_1669_59.ppm}の9画素と，
\code{train_python/week3_6_6.ppm}の上4行24画素を結合した33画素である．
6$\times$6画像の下2行は使用しない．正解値は2枚の色見本画像を最近傍補間で3$\times$3と6$\times$4へ変換し，
画素順を合わせて作成する．100次元以上では各次元の\code{training_data.json}が同じであることを保存済み学習記録と照合した．
評価入力は\code{dimention_40/40_dimention_before.ppm}に由来し，全条件で共通の\code{evaluation_data.h}へ固定した．
この9画素の入力RGBは，上記33画素の学習入力と一致しない．一方，学習と評価は同じ3$\times$3色見本の別測定であり，
未知の色・未知の照明条件を網羅した独立評価ではない．初期40〜70次元については，100次元以降と同程度に詳細な学習実行記録がそろっていないため，
学習条件がすべて同一だったとは断定しない．

\subsection{実験の手順}
\begin{enumerate}[leftmargin=*,itemsep=2pt]
\item 非圧縮の$3\to n\to n\to3$を40，50，60，70次元で確認する．配列長の不整合によるエラーと，RAM配置に起因するリンクエラーを区別する．
\item 70次元の同じ重みから第1・第2層を枝刈りし，COOとCSRを生成する．値・位置情報から枝刈り後の行列を正しく復元できることを検査する．
\item 各圧縮スケッチをビルドし，成功条件では書き込み後に\code{i}で次元数・閾値・容量を確認する．\code{r}で計時とPPM生成を実行し，\code{MEASUREMENT_COMPLETE}を確認する．
\item 100次元で4閾値を比較し，200〜700次元の主系列を閾値0.02で比較する．800・900では保存済みビルドエラーを確認し，700では残りの閾値も検証する．
\item 生のシリアルログ，PPM，ビルド・書き込みログを条件別に保存し，同次元数・同閾値のCOOとCSRを比較する．生成・PC検証だけの条件を実機結果へ含めない．
\end{enumerate}

\subsection{MSE・保存容量・時間の定義}
出力は予測RGBを255倍し，0〜255に制限して小数部を切り捨て，最大値255のP3形式PPMとする．
MSEは9画素のRGB計27成分について

\[
    \operatorname{MSE}=\frac{1}{27}\sum_{p=1}^{9}\sum_{c\in\{R,G,B\}}
    (I_{p,c}-I^{\mathrm{ref}}_{p,c})^2
\]
で求める．計画書の\code{mse_calculator.html}が行う二乗誤差平均と同じ定義で，
報告書作成時は保存PPMから再計算して記録値を照合した．学習時の0〜1のMSEとは尺度と入力が異なる．
保存容量は第1・第2層の値・索引，第3層の重み，全バイアスの\code{sizeof}合計である．
評価入力108 B，画像バッファ27 B，中間配列$8n$ B，プログラムやライブラリの領域はパラメータ容量へ含めない．
静的RAMとフラッシュはビルド結果を別記する．静的RAMだけで実行時ピークRAMを表すことはできない．
\code{micros()}\cite{micros}で，各セットの1,000回の繰り返しの直前・直後を計時する．
経過時間を$u_\ell$とすると，$T_\ell=u_\ell/1000$，平均時間は$\overline{T}=\frac15\sum_{\ell=1}^{5}T_\ell$とする．
計測には入力選択，ループ，\code{predict()}，出力の\code{volatile}変数への集計が含まれる．
センサー読み取り，LED待機，PPM生成，シリアル出力は含まれない．純粋な積和演算だけの時間ではなく，
共通の測定処理を含む1入力当たりの時間である．

\section{実装方法とソースコード}
\subsection{参照したコードの版と役割}
実験プログラムはGitHubの\code{senhari08038010-spec/Agilework2}\cite{repo}を参照した．
ローカルのリモート設定とGitHub上のコミットを確認し，参照版を\code{0446b8663190237bfd36df734d811e808b248883}に固定した．
以下のパスはリポジトリ内の実験\code{/}からの相対パスである．

\begin{table}[H]\centering\small
\caption{主要なソースコードと保存ファイル}
\renewcommand{\arraystretch}{1.2}
\begin{tabularx}{\linewidth}{p{70mm}Y}\toprule
ファイル&役割\\\midrule
\code{train_python/}学習\code{/700_700/}\newline\code{train_L1_normalization_700_700.ipynb}&データを33画素へ整形し，3つのLinearと2つのReLUで学習．重みを行優先で出力\\\midrule
\code{train_python/compress_weights.py}&元ヘッダーの形状検査，閾値処理，COO・CSR生成，索引型選択，復元と容量の検査\\\midrule
\code{scanner_coo_700/scanner_coo_700.ino}&COOの値と行・列番号から直接積和演算\\\midrule
\code{scanner_csr_700/scanner_csr_700.ino}&CSRの行開始位置を用いた行ごとの積和演算\\\midrule
各スケッチの\code{measurement.h}，\code{evaluation_data.h}&共通入力，時間測定，容量集計，PPM生成，有限値と索引の検査\\\midrule
\code{train_python/}学習\code{/}各次元\code{/}\newline\code{compressed_}各閾値\code{/measurements/}&\code{measurement.json}，生ログ，PPM，ビルド結果を条件ごとに保存\\\bottomrule
\end{tabularx}\end{table}

\subsection{変換と疎行列からの直接推論}
\code{compress_weights.py}は元の第1・第2層を行順・列順で走査し，閾値以上かつ非ゼロの値だけを採用する．
同じ走査結果から両形式を作るため，枝刈りによる差と形式による差を混同しない．
元ヘッダーのSHA-256と次元数も生成先に保存する．
整数索引は，行・列番号の最大値と非ゼロ数のそれぞれを表せる最小の符号なし整数型にする．
Arduino側は密行列を復元せず，疎形式を直接参照する．COOの1層分は次の処理に対応する．

\begin{lstlisting}[language=C++]
for (uint32_t r = 0; r < rows; ++r) h[r] = bias[r];
for (uint32_t k = 0; k < nnz; ++k)
    h[row_index[k]] += values[k] * input[col_index[k]];
for (uint32_t r = 0; r < rows; ++r)
    if (h[r] < 0.0f) h[r] = 0.0f;
\end{lstlisting}

CSRでは各行の開始位置と終了位置を取り出し，その範囲だけを計算する．

\begin{lstlisting}[language=C++]
for (uint32_t r = 0; r < rows; ++r) {
    float sum = bias[r];
    for (uint32_t k = row_ptr[r]; k < row_ptr[r + 1]; ++k)
        sum += values[k] * input[col_index[k]];
    h[r] = sum > 0.0f ? sum : 0.0f;
}
\end{lstlisting}

両形式とも第1層と第2層に同じ方式を適用する．第3層は密な配列のまま全要素を計算する．
空行ではバイアスにReLUを適用した値が残るため，非ゼロ重みがないニューロンも無条件には省略しない．
配列長は\code{static_assert}で検査し，次元数は生成ヘッダーから取得する．計時前には索引範囲と，
CSRの開始位置の単調性・終端を検査する．出力のNaN・無限大は異常終了とし，正常な測定完了として記録しない．
元の非圧縮配列は圧縮スケッチに重複して保持しない．

\section{実験結果}
\subsection{非圧縮版の搭載範囲}
\begin{table}[H]\centering\small
\caption{両隠れ層を同じ幅にした非圧縮版の結果}
\begin{tabular}{rrrrrl}\toprule
$n$&容量[B]&MSE&時間[$\mu$s/回]&実機&根拠\\\midrule
40&7,692&427.963&472.436&可&Excel記録・保存PPM\\\midrule
50&11,612&53.593&694.299&可&Excel記録・保存PPM\\\midrule
60&16,332&53.222&958.052&可&Excel記録・保存PPM\\\midrule
70&21,852&---&---&未実施&RAM不足のリンクエラー\\\bottomrule
\end{tabular}\end{table}

40〜60次元のMSEは保存PPMから再計算し，時間はExcelの記録値を用いた．
元の5セット分の生ログは確認できないため，初期条件の時間のばらつきは算出しない．
70次元は非圧縮でビルドできず，過去のエラー出力では\code{.heap}と\code{.stack_dummy}の重複が確認される．
単純な配列サイズ不一致のエラーとは区別した．試験した40，50，60，70の範囲では，最後に動作した非圧縮モデルは60次元である．

\subsection{閾値0.02での次元数比較}
\begin{table}[H]\centering\small
\caption{保存済み実機ログから照合した圧縮版の結果（$t=0.02$）}\label{tab:main}
\setlength{\tabcolsep}{5pt}
\begin{tabular}{rrrrrr}\toprule
$n$&形式&容量[B]&圧縮率[\%]&MSE&時間[$\mu$s/回]\\\midrule
100 & COO & 2,744 & 6.350 & 83.370 & 192.333\\\midrule
100 & CSR & 2,824 & 6.535 & 83.370 & 206.023\\\midrule
200 & COO & 4,738 & 2.847 & 94.889 & 336.139\\\midrule
200 & CSR & 5,019 & 3.016 & 94.889 & 370.773\\\midrule
300 & COO & 7,029 & 1.902 & 50.593 & 488.635\\\midrule
300 & CSR & 7,365 & 1.993 & 50.593 & 531.583\\\midrule
400 & COO & 9,078 & 1.391 & 62.407 & 635.957\\\midrule
400 & CSR & 9,598 & 1.470 & 62.407 & 697.766\\\midrule
500 & COO & 10,832 & 1.066 & 53.852 & 767.470\\\midrule
500 & CSR & 11,616 & 1.143 & 53.852 & 851.307\\\midrule
600 & COO & 12,911 & 0.885 & 103.185 & 915.090\\\midrule
600 & CSR & 13,877 & 0.951 & 103.185 & 1016.239\\\midrule
700 & COO & 14,813 & 0.747 & 38.037 & 1054.399\\\midrule
700 & CSR & 16,003 & 0.807 & 38.037 & 1175.937\\\midrule

\bottomrule\end{tabular}
\end{table}
表\ref{tab:main}の全条件はビルド・書き込み・実機測定が完了しており，
各条件のPPM，\code{MEASUREMENT_COMPLETE}を含むシリアルログと\code{measurement.json}が保存されている．
今回，PPMからのMSE，5セットの経過時間からの平均，生成配列からの容量を照合し，記録値と一致した．同じ次元数・閾値では，
COOとCSRの27成分はすべて一致した．
MSEは100次元で83.370，200で94.889，300で50.593となり，次元数に対して単調に減少しなかった．
比較した条件のうち700次元・閾値0.02の38.037が最小だった．一方，推論時間と保存容量はこの主系列では次元数とともに増加した．

\begin{figure}[H]\centering
\begin{tikzpicture}
\begin{axis}[width=.88\linewidth,height=55mm,xlabel={両隠れ層の次元数 $n$},ylabel={MSE},xmin=80,xmax=720,ymin=0,ymax=120,xtick={100,200,300,400,500,600,700},grid=major]
\addplot[black,mark=*,thick] coordinates {(100,83.37037037037037) (200,94.88888888888889) (300,50.592592592592595) (400,62.407407407407405) (500,53.851851851851855) (600,103.18518518518519) (700,38.03703703703704)};
\end{axis}\end{tikzpicture}
\caption{閾値0.02の次元数とMSE．COO・CSRは同値のため1本で示す．}
\end{figure}
\begin{figure}[H]\centering
\begin{tikzpicture}
\begin{axis}[width=.88\linewidth,height=55mm,xlabel={両隠れ層の次元数 $n$},ylabel={平均推論時間 [$\mu$s/回]},xmin=80,xmax=720,ymin=0,xtick={100,200,300,400,500,600,700},grid=major,legend pos=north west]
\addplot[black,mark=*,thick] coordinates {(100,192.3334) (200,336.1394) (300,488.635) (400,635.957) (500,767.4702) (600,915.0898) (700,1054.3992)};\addlegendentry{COO}
\addplot[black,dashed,mark=square*,thick] coordinates {(100,206.023) (200,370.773) (300,531.5828) (400,697.7656) (500,851.307) (600,1016.2392) (700,1175.9366)};\addlegendentry{CSR}
\end{axis}\end{tikzpicture}
\caption{閾値0.02の次元数と平均推論時間．}
\end{figure}

\subsection{閾値による変化}
\begin{table}[H]\centering\small
\caption{100・700次元での閾値比較．容量はCOO／CSRの順．}\label{tab:threshold}
\setlength{\tabcolsep}{4pt}
\begin{tabular}{rrrrrr}\toprule
$n$&$t$&容量[B]&共通MSE&COO時間&CSR時間\\\midrule
100 & 0.005 & 2,984 / 3,024 & 91.778 & 206.539 & 217.722\\\midrule
100 & 0.01 & 2,780 / 2,854 & 91.778 & 194.463 & 207.778\\\midrule
100 & 0.02 & 2,744 / 2,824 & 83.370 & 192.333 & 206.023\\\midrule
100 & 0.03 & 2,702 / 2,789 & 88.481 & 189.779 & 203.934\\\midrule
700 & 0.005 & 35,632 / 32,317 & 未測定 & 未測定 & 未測定\\\midrule
700 & 0.01 & 14,861 / 16,039 & 38.481 & 1056.789 & 1177.569\\\midrule
700 & 0.02 & 14,813 / 16,003 & 38.037 & 1054.399 & 1175.937\\\midrule
700 & 0.03 & 14,783 / 15,981 & 42.852 & 1052.777 & 1174.722\\\midrule

\bottomrule\end{tabular}\par\smallskip
時間の単位は$\mu$s/回．700・0.005の容量は生成配列からの計算値．
\end{table}

100次元では0.005と0.01のMSEは同じ91.778であり，0.02で83.370，0.03で88.481となった．
700次元でも0.01から0.02でMSEがわずかに小さくなった後，0.03では42.852へ増加した．
閾値の増加は容量を減らすが，MSEを常に悪化させるわけでも，常に改善させるわけでもなかった．
70次元については，Excelに0.005，0.01，0.02，0.03のMSEがそれぞれ83.07，83.07，84.41，90.30と記録されている．
この系列の容量はCOOが2,186，2,120，2,102，2,078 B，CSRが2,199，2,144，2,129，2,109 Bである．
ただし0.005の実機PPM・生ログは未確認であり，この値は測定記録として区別する．
70次元の0.02から100次元の0.02へのMSE差だけを，モデルの幅による効果と解釈することは避ける．

\subsection{RAM使用量と測定できなかった条件}
\begin{table}[H]\centering\small
\caption{閾値0.02のビルド時メモリ量．対の値はCOO／CSRの順．}
\setlength{\tabcolsep}{5pt}
\begin{tabular}{rrrr}\toprule
$n$&静的RAM[B]&中間配列[B]&フラッシュ[B]\\\midrule
100 & 9,516 / 9,600 & 800 & 58,448 / 58,564\\\midrule
200 & 11,512 / 11,792 & 1,600 & 60,460 / 60,756\\\midrule
300 & 13,804 / 14,140 & 2,400 & 62,768 / 63,120\\\midrule
400 & 15,852 / 16,372 & 3,200 & 64,824 / 65,352\\\midrule
500 & 17,604 / 18,392 & 4,000 & 66,568 / 67,372\\\midrule
600 & 19,684 / 20,652 & 4,800 & 68,648 / 69,632\\\midrule
700 & 21,588 / 22,776 & 5,600 & 70,552 / 71,756\\\midrule

\bottomrule\end{tabular}
\end{table}

700次元・0.02の中間配列は$700\times2\times4=5,600$ Bである．
静的RAMにこの配列だけを加えるとCOOは27,188 B，CSRは28,376 Bであり，
32,768 Bまでの差はそれぞれ5,580 B，4,392 Bとなる．この差には他のスタック・ヒープ領域も収める必要があり，
実測の空きRAMやピーク使用量を示すものではない．
800次元・0.02と900次元・0.02はCOOのリンク段階でRAM領域の不足・重複エラーが発生し，
実機への書き込み前に停止した．生成容量は800でCOO 16,941 B／CSR 18,295 B，900で19,067 B／20,589 Bである．
CSRの同条件の実機試験は実施されていないため，CSRもリンク失敗したとは記載しない．
700次元・0.005はCOOのパラメータ容量だけで35,632 Bとなり，SRAMの32,768 Bを超える．
CSRの生成容量32,317 Bも，中間配列5,600 Bだけを加えると37,917 Bとなるが，
CSRについては実際のビルド・実機失敗を確認した結果ではなく，RAM配置を維持した場合の容量上の見積もりである．
この条件もCOOの失敗時点で測定を停止している．
実機で確認できた最大の幅は両形式とも700であり，非圧縮で確認できた60に対して各隠れ層が640増えた．
これは今回の条件で確認できた範囲を表し，形式固有の絶対的な最大値ではない．

\section{考察}
\subsection{メモリ縮減の手計算と搭載範囲}
700次元・閾値0.02の全パラメータ容量を，実装の型と要素数から計算する．枝刈り後の非ゼロ数は第1層47，
第2層59である．第1層の行番号と第2層の行・列番号は2 B，第1層の列番号は1 B，
両層のCSR開始位置は非ゼロ数が255以下のため1 Bとなる．非圧縮で残す第3層と全バイアスは

\[
    4(3\times700+700+700+3)=14,012\ \mathrm{B}
\]
である．したがって

\begin{align*}
    S_{\mathrm{COO}}&=47(4+2+1)+59(4+2+2)+14,012\\
    &=329+472+14,012=14,813\ \mathrm{B},\\
    S_{\mathrm{CSR}}&=47(4+1)+701+59(4+2)+701+14,012\\
    &=235+701+354+701+14,012=16,003\ \mathrm{B}.
\end{align*}
非圧縮容量は$4(700^2+8\times700+3)=1,982,412$ Bなので，
圧縮率はCOOが0.747\%，CSRが0.807\%，削減率は99.253\%，99.193\%となる．
いずれも同じ700次元の非圧縮モデルを分母とした理論容量との比較であり，
非圧縮700次元を実機で動作させた比較ではない．
第1層と第2層のゼロ割合はそれぞれ97.762\%，99.988\%と高い．幅は大きくても実際に残る接続は少なく，
「大きな密ネットワークと同等の計算量を搭載した」と解釈してはならない．
また，COOでは残した第3層と全バイアスが14,813 B中14,012 B，約94.6\%を占める．
以後の小容量化では，既に小さい疎配列だけを削るより，この残存部分と中間配列に着目する必要がある．

\subsection{COOが小容量・短時間だった理由}
700次元ではCSRの開始位置配列が$701\times2=1,402$ B必要だが，COOの行番号は$(47+59)\times2=212$ Bで済む．
他の値・列番号が共通なので，差は$1,402-212=1,190$ Bとなり，実測した16,003 Bと14,813 Bの差を説明できる．
このモデルでは空行が多いため，CSRが行番号を省略する利点よりも，全行の開始位置を保存するコストが大きい．
同条件の平均推論時間はCOO 1,054.399 $\mu$s，CSR 1,175.937 $\mu$sであり，COOはCSR比で約10.34\%短い．
容量もCSR比で約7.44\%小さい．CSRは各行で開始位置と終了位置を参照するため，
非ゼロ要素が少ない条件では行管理の処理が相対的に大きくなると考えられる．ただし，
時間差の原因を特定する命令単位の計測はしていない．これは本実装と高い疎性の下での説明であり，
すべての疎行列でCOOが速いという結論にはならない．
700次元・0.02の5セットの時間は，COOが1054.401, 1054.398, 1054.401, 1054.397, 1054.399，
CSRが1175.935, 1175.937, 1175.937, 1175.937, 1175.937 $\mu$s/回であった．
セット間の標本標準偏差はそれぞれ0.0018，0.0009 $\mu$s/回である．同一書き込み・同一測定内の反復であり，
別日に再測定した再現性や学習し直した際のばらつきを示すものではない．

\subsection{MSEの一致と次元数に対する非単調性}
同じ閾値のCOO・CSRで出力が一致したことは，両形式が同じ枝刈り後の重みを表す設計と整合する．
これは保存形式間の整合性を示すが，枝刈り前の品質が保持されたことを直接示すものではない．
700次元の枝刈り前モデルをRAMへ搭載できず，同モデルの枝刈り前後を実機で比較していないためである．
次元数を増やしてもMSEが単調に改善しない理由として，幅ごとに最適化経路や得られる重みが異なること，
固定のL1係数でも重み数と重みの総和が変わること，閾値0.02が各モデルに与える影響が異なることが考えられる．
seedを同じ0にしても，形状が違うモデルに同一の初期重みが与えられるわけではない．
これらを切り分ける比較は実施していないため，300次元での改善や600次元での悪化の原因は仮説として扱う．
同じモデルの0.005と0.01で整数画像のMSEが一致する場合も，枝刈りした接続が評価入力に対して寄与しないことや，
浮動小数点の出力差が整数化で消えることが候補となる．整数PPMの一致だけから内部出力まで完全一致したとは限らない．
閾値を変更した際はモデル識別情報と非ゼロ数も確認する必要がある．

\subsection{評価の範囲と限界}
精度評価は同じ9画素と色見本に限定されており，複数画像・照明・測定日の条件を変えた一般化性能は評価していない．
また，各幅で学習したモデルは1個であり，5セットの時間測定は学習を5回繰り返した実験ではない．
700・0.02が最良という結論も，今回比較した幅・閾値・9画素の範囲に限られる．
初期のExcelには，第1隠れ層40固定の系列と両層を変更した系列が併記されている．本報告では両層系列を選び，
圧縮率の小数表記を百分率へ換算した．40固定・50次元の「保存容量」欄64,620 Bは過去のフラッシュ使用量の発言と一致し，
理論パラメータ容量9,452 Bと異なるため，パラメータ容量の比較には採用していない．
記録の名称と集計範囲を確かめずに異なる系列を混在させると，圧縮効果を誤評価するおそれがある．

\section{さらなる性能向上に向けた検討}
\begin{enumerate}[leftmargin=*,itemsep=4pt]
\item 枝刈り前後の精度差を分離する．同じモデルをPC上の同一float32演算で比較し，Arduinoに搭載可能な幅では実機でも確認する．整数PPMのMSEに加え，整数化前の出力差を測る．枝刈り後の再学習も別条件として検証する．
\item 評価入力を拡張する．学習と異なる色や照明，別の色見本，複数回の実測を用い，複数seedで再学習する．各条件のMSEの平均とばらつきを示し，次元数と最適化の影響を分ける．
\item 残る14,012 Bの削減を検討する．第3層の枝刈り・量子化，隠れニューロン単位の構造化枝刈りを比較する．ニューロン削除はバイアスや次層への影響があるため，単に入力重みが0という理由だけでは削除しない．
\item Flashへの重み配置や8／16 bit量子化を別の実験条件として調べる．RAM削減と引き換えに読出し時間や数値誤差が変わるため，今回のRAM配置の結果とは分けて評価する．中間配列・スタックのピーク量も測定する．
\item 限界付近を細かく検証する．700〜800の間を10刻み等で調べ，CSRも独立にビルド・実機確認する．各形式の最大動作幅は，終了条件をそろえた試験後に比較する．
\end{enumerate}

\section{結論}
両隠れ層を同じ幅に拡大するRGB補正モデルに対し，第1・第2層の枝刈りとCOO・CSR保存を実装した．
試験範囲では非圧縮60次元に対し，圧縮版は両形式とも700次元の実機動作を確認した．700次元・閾値0.02では，
理論非圧縮容量1,982,412 Bに対してCOOは14,813 B，CSRは16,003 Bとなり，MSEはともに38.037であった．
COOはCSRより容量が約7.44\%小さく，平均推論時間が約10.34\%短かった．
この差は，非ゼロ要素が非常に少ないモデルでは，CSRの全行分の開始位置配列がCOOの行番号配列より大きくなることと整合した．
一方，次元数とMSEは単調な関係を示さず，枝刈り前後の品質保持と形式ごとの絶対的な搭載限界は本実験だけでは確定できない．
今後は評価データと学習反復を増やし，残存パラメータと作業領域の削減を検討する必要がある．

\begin{thebibliography}{9}\small
\bibitem{plan} 泉水悠斗：アジャイルワーク２ 実験計画書，実験計画書\code{/}実験計画書\code{.tex}，参照版は文献\cite{repo}のコミットに含まれる．
\bibitem{han} Han, S., Pool, J., Tran, J. and Dally, W. J.: Learning both Weights and Connections for Efficient Neural Networks, Advances in Neural Information Processing Systems 28 (2015), \url{https://arxiv.org/abs/1506.02626}（参照2026-11-13）．
\bibitem{coo} SciPy: coo\_matrix, SciPy Manual, \url{https://docs.scipy.org/doc/scipy/reference/generated/scipy.sparse.coo_matrix.html}（参照2026-11-13）．
\bibitem{csr} SciPy: csr\_matrix, SciPy Manual, \url{https://docs.scipy.org/doc/scipy/reference/generated/scipy.sparse.csr_matrix.html}（参照2026-11-13）．
\bibitem{board} Arduino: UNO R4 WiFi, \url{https://docs.arduino.cc/hardware/uno-r4-wifi/}（参照2026-11-13）．
\bibitem{micros} Arduino: micros(), Language Reference, \url{https://docs.arduino.cc/language-reference/en/functions/time/micros/}（参照2026-11-13）．
\bibitem{repo} 泉水悠斗：Agilework2，実験ソースコード・測定記録，\url{https://github.com/senhari08038010-spec/Agilework2/tree/0446b8663190237bfd36df734d811e808b248883}（非公開リポジトリ，参照2026-11-13）．
\end{thebibliography}

\clearpage
\appendix
\section{結果の参照先と照合方法}
初期の非圧縮・70次元の記録は，Codex作業ディレクトリに保存された\code{outputs/01a16723-9148-7903-bf48-1dcddd5a209e/}実験結果\code{.xlsx}を参照した．両隠れ層系列は同シートの12〜23行である．
非圧縮40，50，60のPPMは，リポジトリ内実験\code{/train_python/}学習\code{/40_40/}等の保存画像と照合した．
70次元の非圧縮PPMファイルは有効な画素を含まないため，実機成功の証拠には用いていない．
100次元以上の数値の一次資料は，文献\cite{repo}の実験\code{/train_python/}学習\code{/}以下にある条件別の記録である．
例えば700次元・0.02は，\code{700_700/compressed_t1_0p02_t2_0p02/measurements/2026-11-06_165736_JST/}に
COO・CSR両方のPPM，\code{serial_measurement.txt}，\code{measurement.json}，\code{arduino_build.txt}がある．
100次元の4閾値，200〜600の0.02，700の0.01・0.02・0.03の計24件を照合した．
モデルの元ヘッダー，枝刈り後の層別非ゼロ数・索引型と容量は各条件の\code{compression_report.json}，
PC上の復元・推論照合は\code{validation_report.json}，実機失敗条件は\code{memory_error.json}と
生のビルドログを参照した．これらの生成・PC検証ファイルにある「実機未測定」のフラグは，
後から保存した実機測定ログとは役割が異なる．成功の判定には実機側の記録を使用した．今回の再計算値と参照パスは，
報告書に付属する照合記録\code{.json}に保存した．
\end{document}
